% Section 9.3: A dictionary of conventions | Part: 1 | Depends on: sec:vectors to sec:duality | Status: done | Original lines: 1753-1781
%!TEX root = ../../../main.tex
\subsection{A dictionary of conventions}\label{sec:conventions}

Because the three sources were written independently, the same object appears with different normalisations. They have been unified
as follows; this is the place to look when a formula elsewhere in the literature seems to disagree by a factor or a sign.

% REORDERED: Remark rem:kfactorial now precedes the table that refers to it.
\begin{remark}[the factor $k!$ between the two realisations of $\bigwedge$]\label{rem:kfactorial}
Let $\alpha^1,\dots,\alpha^k\in V^*$ and $v_1,\dots,v_k\in V$. The form convention of Section~\ref{sec:ext} gives $(\alpha^1\wedge\cdots\wedge\alpha^k)(v_1,\dots,v_k)=\det[\alpha^i(v_j)]$. If instead one embeds multivectors in $T^k(V)$ by
$v_1\wedge\cdots\wedge v_k\mapsto\frac1{k!}\sum_\sigma\sgn(\sigma)\,v_{\sigma(1)}\otimes\cdots\otimes v_{\sigma(k)}$ and evaluates the resulting tensor on $(\alpha^1,\dots,\alpha^k)$, one obtains $\frac1{k!}\det[\alpha^i(v_j)]$. The two realisations are identified by multiplying degree-$k$ elements by $k!$: this map is an algebra isomorphism, because the factor $\frac{(k+l)!}{k!\,l!}$ in \eqref{eq:altwedge} satisfies $(k+l)!=k!\,l!\cdot\frac{(k+l)!}{k!\,l!}$. Nothing in the abstract algebra changes; only the numbers one gets by evaluating on tensors do.
\end{remark}

% AUTHOR-NOTE: the signature row says "GR notes: (negative, positive)", but the formal box of Section sec:metric
% (original L794) uses (positive, negative) for the GR notes. Check which one the GR notes use (docs/AUTHOR_NOTES.md #3).
\begin{center}
\small
\renewcommand{\arraystretch}{1.3}
\begin{tabular}{@{}p{3.1cm}p{5.4cm}p{5.8cm}@{}}
\toprule
\textbf{Item} & \textbf{Convention in this document} & \textbf{Other conventions met in the sources}\\
\midrule
Signature $(p,q)$ & $p$ positive, $q$ negative; Euclidean $(n,0)$ (pitfall in Section~\ref{sec:quadform}) & GR notes: (negative, positive); mostly minus $(1,3)$ in the spacetime algebra\\
Wedge, forms & $\dd x\wedge\dd y=\dd x\otimes\dd y-\dd y\otimes\dd x$; $\det$ pairing & Kobayashi--Nomizu: factor $\frac12$\\
Wedge, multivectors & $v\wedge w\mapsto\frac12(v\otimes w-w\otimes v)$ in $T(V)$; $a\wedge b=\frac12(ab-ba)$ & see Remark~\ref{rem:kfactorial}\\
Tensor type $(\ell,k)$ & $\ell$ contravariant (up), $k$ covariant (down) & some texts reverse the order\\
Dual basis & $e^i$ ($=\dd x^i$ in coordinates) & $e^{*i}$ in the GR notes\\
Hodge star vs.\ dual & $\hodge A=\rev A\pss=(-1)^{k(k-1)/2+n(n-1)/2}\dual A$ & dual with $\pss$ or $\pss^{-1}$ on either side\\
Curvature & $R(X,Y)Z=\nabla_X\nabla_YZ-\nabla_Y\nabla_XZ-\nabla_{[X,Y]}Z$ & Carroll/MTW~\cite{Carroll,MTW} agree; others differ in index order (Section~\ref{sec:curv})\\
$\nabla$ & plain: covariant derivative. Bold $\vd$: vector derivative & Hestenes~\cite{HS} writes $\nabla$ for $\vd$\\
\bottomrule
\end{tabular}
\end{center}
